\section{Exact image and fiber structure}

The cubic discriminant is $-4Q$, where
\[
Q=27a^2c^2-18abc+16a+b^3c-b^2.
\]
The triple-root locus is
\[
\Gamma=V(3bc-4,\,12a-b^2).
\]

\begin{theorem}\label{thm:image}
Over $\mathbb C$,
\[
F(\mathbb C^3)=\mathbb C^3\setminus\Gamma,
\qquad
S_F=V(Q),
\]
where $S_F$ is the nonproperness set.  Fibers have respectively three, one,
and zero points on $Q\ne0$, on $V(Q)\setminus\Gamma$, and on $\Gamma$.
\end{theorem}

This follows from Theorem~\ref{thm:marked-root}: a nonzero binary cubic has
three, one, or zero simple roots according to multiplicity type $(1,1,1)$,
$(2,1)$, or $(3)$.  Repeated roots are precisely the poles of the affine
reconstruction chart and give escaping branches.  Away from the discriminant,
all roots remain simple, while the projective infinity root is controlled by
the second regular chart.
